\sscp{Some implications of **mantəR ‘cuscus’}{14-04-05}
The \tsc{\ili{Muna}-Buton} word \it{mante} ‘bear cuscus’ in seven languages
directly undermine \citeauthor{Blust2012b}’s (\citeyear[274]{Blust2012b})
criticism that \citeauthor{Schapper2011}’s (\citeyear{Schapper2011}) scenario about marsupial terms,
“implies an entry via Sulawesi into the Lesser Sundas,
with subsequent eastward movement into the Moluccas.
If this were the case,
we would expect some linguistic trace in Sulawesi,
but there is none; despite the speculations in \citet{DonohueGrimes2008}
about possible subgrouping ties between at l\ili{east} some of the
languages of Sulawesi and languages in eastern Indonesia,
closer scrutiny of the evidence reveals nothing of substance
to support this idea.” We have provided evidence from the historical
phonology in \crf{ch:SB} of links from \ili{Central} Sulawesi-\ili{Banggai}-\ili{Taliabu},
the latter in Maluku \ili{east} of the Blust line.
In this chapter we provide additional lexical evidence through
the distribution of ubiquitous reflexes of **mantəR ‘cuscus’
and closely related forms in \tsc{\ili{Muna}-Buton}
and throughout our Wallacean subgroups (minus \tsc{\ili{Bima}-Lembata}),
as well as \ili{Buru} \it{tonal} ‘cuscus’ borrowed from a south\ili{east} Sulawesi language.\footnote{
		One other borrowing in \ili{Buru} from south\ili{east} Sulawesi is \it{bara}
		`don't, \tsc{prohibative}'
		from \ili{Tukang Besi} \it{bara} `don't, \tsc{prohibative}'.}

Furthermore, the medial *ns
and final *r of **mans(aə)r are not well supported by the fuller
data, whereas **mantəR ‘cuscus’ and close derivatives account for both
Sulawesi and Wallacean subgroups by regular sound correspondences.
The possibly derivative form of **mandəR is explainable from
post-nasal voicing of *nt; **məntəR is explainable by sporadic
vowel shift in one or two geographic regions.
The Sulawesi forms preserve a medial NC cluster in \it{mante},
and in that regard appear to be more conservative than the forms
in \tsc{Sula-Buru} and \tsc{\ili{Ambon}-Seram} languages.
The Sulawesi forms also require *əR{\#} to derive *ə > \it{e},
instead of the otherwise expected *ə > \it{o}.
Many \tsc{\ili{Ambon}-Seram} and \tsc{Seram-Tanimbar-Bomberai} languages
regularly preserve the final *R{\#} that is regularly lost in \tsc{\ili{Muna}-Buton} languages.

The gloss of ‘bandicoot’ for putative PCEMP **mans(aə)r must
be rejected from the evidence in Wallacean subgroups
and Sulawesi.
There is overwhelming support for the gloss of ‘cuscus’ from
eight subgroups in Sulawesi
and Wallacea (see Tables \ref{tab:14-07a},
\ref{tab:14-07b}, \ref{tab:14-07c}, and \ref{tab:14-11}).

We do not dispute the gloss of the \tsc{POc} etymon of *mʷajar as ‘bandicoot’.
(Note that \tsc{POc} *j is a prenasalised palatal stop [ɲʤ].)
%hence /ᶮʤ/ → [ɲʤ].)
%hence /\su{ɲ}ʤ/ → [ɲʤ].)
However, methodologically we do not We do not accept
that the gloss of \tsc{POc} *mʷajar ‘bandicoot’ should
determine the gloss for **mantəR,
which has overwhelming support for ‘cuscus’.
%and only one highly irregular form outside of Oceanic languages
%that is glossed ‘bandicoot’.

It appears that the medial *ns of putative PCEMP **mans(aə)r,
along with its gloss of ‘bandicoot’ has been determined by \tsc{POc} *mʷajar ‘bandicoot’.
\citet[27--46]{ RossPawleyOsmond2023} is the most recent discussion
of \tsc{POc} proto-phonology and its antecedents.
He shows that \tsc{POc} *j is a reflex of PMP *j,
*z, *ns, or *nz, but not normally of PMP *nt or *nd.\footnote{
		The normal reflexes of PMP *nt in \tsc{POc} are *t or *d (three examples each).
		The reflexes of PMP *nd in \tsc{POc} are *r (one example) or
		*dr [ⁿr] (two examples).}
Thus, PMP medial *ns in **mans(aə)r (or perhaps **mansəR) can account
for \tsc{POc} *mʷajar regularly,
while medial *nt cannot.\footnote{
  Malcolm Ross (pers. comm. September 2025) now finds that the
	reconstructable form for \tsc{POc} ‘bandicoot’ is *mʷajaŋ with final *ŋ.
	This form also find support in EMP from \ili{Ambai}
	\it{mansaŋ} {\tl} \it{mansa} ‘grizzled tree kangaroo, \it{Dendrolagus inustus}’.
	While medial *ns can account for the medial consonant
	of \tsc{POc} *mʷajaŋ, final *R cannot account for the final consonant.
	If and how \tsc{POc} *mʷajaŋ `bandicoot' is connected with 
	**mantəR ‘cuscus’ remains to be seen.}

However, determining the gloss
and form for the languages of Wallacea on the basis
of the Oceanic data alone is methodologically unsound.
In support of medial *ns we could add \ili{Ambai} (SHWNG) \it{mansaŋ
{\tl} mansa} ‘grizzled tree kangaroo, \it{Dendrolagus inustus}’ \citep{Price2002}.
So medial *ns is supported in EMP,
but not in Wallacean subgroups, nor in \tsc{\ili{Muna}-Buton}.
Even with the \ili{Ambai} data,
the gloss of ‘bandicoot’ is not supported outside of Oceanic languages.
Collectively EMP together would be only one witness.
From the perspective of the \tsc{\ili{Muna}-Buton} languages
and the Wallacean subgroups,
there is evidence of a shift in form
and meaning from \tsc{\ili{Muna}-Buton} **mantəR ‘bear cuscus (specific)’
> Wallacean subgroups **mantəR (and similar forms) ‘cuscus (generic)’,
broadening in meaning > EMP **mansəR/**mansaR ‘marsupial (more
generic)’, and then a further shift in form
and narrowing in meaning > \tsc{POc} *mʷajor/*mʷajar ‘bandicoot
(specific)’.\footnote{
		Malcolm Ross (pers. comm. September 2025) now finds that
		the reconstructable form for \tsc{POc} ‘bandicoot’ is *mʷajaŋ.
		Wallacean **mantəR ‘cuscus’ is different
		in both form and meaning from \tsc{POc} *mʷajaŋ ‘bandicoot’.
		Yes, they are both marsupials,
		but one lives in the trees,
		and the other on the ground.
		If the two etyma are related,
		there has been a significant shift in both form
		and meaning, and a credible explanation of the pathway from medial
		**nt > \tsc{POc} *j and final **R > \tsc{POc} *ŋ would be needed.}

Where showing continuity is possible,
it is also desirable.
But the fuller data do not permit that in this case.
Of course, the earlier reconstructions pertaining to marsupial
terms were not showing awareness of the Sulawesi data,
nor of the fuller data
and corrected glosses of the Wallacean forms that we present
in this chapter and in appendix \ref{ch:SurMar}.